\section{A bounded residue coupling for log concave weights}\label{spb:ml:sec:coupling}
This appendix gives a general supplementary lemma. It is not needed for the parity specific proof above, but clarifies why very rare residue conditioning can still preserve coarse location and scale. Its ingredients are classical log concavity and monotone likelihood ratio order; see \cite{sw,dm}. We give the full modular argument, without a novelty claim for this formulation.

\begin{theorem}[Bounded coupling of all present residues]\label{spb:ml:thm:coupling}
Let $w_0,\ldots,w_N$ be strictly positive and log concave, meaning $w_k^2\ge w_{k-1}w_{k+1}$ whenever $1\le k<N$. Normalize them to a probability mass function $X$. Fix an integer $q\ge2$, and let $X_a$ have its conditional law on residue $a$ modulo $q$ whenever that residue is present. There is a simultaneous coupling with
\[
 |X_a-X_b|\le q-1\quad\hbox{almost surely for every present pair }a,b.
\]
For each fixed present residue there is also a coupling with the original law such that $|X_a-X|\le q-1$ almost surely. The common quantile coupling has this property. The statement extends to any translated consecutive finite support.
\end{theorem}
\begin{proof}
We first record a cross product proof of the required stochastic order, including zero endpoints. If probability masses $P,Q$ on the integers satisfy
\[
 P(i)Q(j)\ge Q(i)P(j)\qquad(i<j),
\]
then for every integer $t$,
\[
 F_P(t)-F_Q(t)=\sum_{i\le t<j}\{P(i)Q(j)-Q(i)P(j)\}\ge0.
\]
Thus $P$ is stochastically no larger than $Q$. A decreasing likelihood ratio on overlapping interval supports implies the cross product condition when the numerator support begins and ends no later than the denominator support. This formulation avoids division by zero.

For present residues $0\le a<b\le q-1$, write $X_a=a+qK_a$, $X_b=b+qK_b$ and
$A_a=\lfloor(N-a)/q\rfloor$, $A_b=\lfloor(N-b)/q\rfloor$. The two index supports are $[0,A_a]$, $[0,A_b]$, with $A_a-1\le A_b\le A_a$. Log concavity makes $r_j=w_{j+1}/w_j$ nonincreasing. Therefore for each positive shift $d$, the ratio
$w_{t+d}/w_t=\prod_{j=0}^{d-1}r_{t+j}$ is nonincreasing wherever defined.

On the common support the mass ratio of $K_b$ to $K_a$ is a positive constant times $w_{b+qk}/w_{a+qk}$, hence decreasing. Both supports start at zero, and the numerator ends no later. The cross product criterion gives $K_b\le_{\mathrm{st}}K_a$.

Next compare $K_a-1$, supported on $[-1,A_a-1]$, with $K_b$, supported on $[0,A_b]$. The former begins and ends no later. On their overlap its likelihood ratio is a constant times
\[
 \frac{w_{a+q(k+1)}}{w_{b+qk}},
\]
whose shift $q+a-b$ lies in $\{1,\ldots,q-1\}$; this ratio is decreasing too. At the left endpoint only the first distribution may have mass, and at a possible unmatched right endpoint only the second may have mass. Thus all endpoint cross products have the same orientation, and
\begin{equation}\label{spb:ml:eq:interlacing}
 K_a-1\le_{\mathrm{st}}K_b\le_{\mathrm{st}}K_a.
\end{equation}
This includes singleton supports and empty overlaps. Residues with no support are simply omitted.

Let $Q_a(u)$ be the lower quantile of $K_a$ for $0<u<1$. Stochastic order is equivalent to pointwise quantile order, so \eqref{spb:ml:eq:interlacing} gives
$Q_a(u)-1\le Q_b(u)\le Q_a(u)$. With a single uniform random variable $U$, set $X_a=a+qQ_a(U)$ for every present residue. For $a<b$, their difference is either $b-a$ or $b-a-q$, both of absolute value at most $q-1$. This proves the simultaneous bound.

Choose independently a residue $J$ with probabilities $\PP(X\equiv j\pmod q)$, and set $X=X_J$ on this coupling. Its law is the original mixture and it lies within $q-1$ of every $X_a$. The ordinary common quantile coupling also works: at each $u$, the mixture quantile lies between the smallest and largest component quantiles. Below the smallest all component distribution functions are less than $u$, and at the largest all are at least $u$. These extreme quantiles differ by at most $q-1$, proving the assertion. Translation changes only residue labels and leaves all distances unchanged.
\end{proof}

\begin{corollary}[Coarse binomial bounds for any fixed modulus]\label{spb:ml:cor:coupling}
For any defined residue conditional of $X\sim\Bin(M,\pi)$, with ordinary mean $\nu$ and variance $v$,
\[
 |\E X_a-\nu|\le q-1,\qquad
 \Var X_a\le2v+2(q-1)^2.
\]
For every interval $I$ of length $L$, there is a constant $C_q$ such that
\[
 \PP(X_a\in I)\le C_q(L+1)/\sqrt{v+1}.
\]
If $q$ is fixed and $v\to\infty$, then $(X_a-\nu)/\sqrt v\Longrightarrow N(0,1)$ for any sequence of present residues.
\end{corollary}
\begin{proof}
Interior binomial weights are positive and log concave because
$w_{k+1}/w_k=\pi(M-k)/((1-\pi)(k+1))$ decreases in $k$. Endpoint probabilities give point masses and need no conditioning argument. In the bounded coupling, the mean bound follows immediately and Minkowski's inequality gives
$\E(X_a-\nu)^2\le(\sqrt v+q-1)^2\le2v+2(q-1)^2$. Variance is no larger. If $X_a\in I$, the coupled $X$ lies in the interval expanded by $q-1$ at each end. The ordinary Fourier atom bound proved in Lemma~\ref{spb:ml:lem:rare} bounds its mass by $C_q(L+1)/\sqrt{v+1}$. Finally the standardized coupling displacement is at most $(q-1)/\sqrt v$, so the ordinary binomial characteristic function central limit argument transfers the limit.
\end{proof}

For $q=2$, these estimates give an alternative to the first part of Section~\ref{spb:ml:sec:global}, with no conditioning probability in any denominator. They do not supply the model specific cubic density expansion. Consequently this appendix makes no global approximation claim for another endpoint parameter or growing modulus. The uniform coupling distance $q-1$ itself cannot be reduced: take $N=q-1$ and a binomial success probability tending to zero. Its residue $q-1$ conditional is concentrated at $q-1$, while the unconditional mean tends to zero.
